\section{Hypotheses and study design}
\subsection{Expected effects and proposed mechanism}

\begin{frame}{Possible changes in model understanding}
\centering

\begin{tikzpicture}[
    profile/.style={
        align=center,
        text width=3.5cm,
        font=\small
    },
    pathway/.style={
        -{Stealth[length=3mm]},
        very thick,
        draw=LeipzigBlue
    },
    explanation/.style={
        align=center,
        text=LeipzigBlue,
        font=\scriptsize\bfseries,
        fill=white,
        inner sep=3pt
    }
]

% Same dimensions and profiles as the earlier slide
\draw[-{Stealth},thick]
    (0,0) -- (11,0) node[right] {FF};
\draw[-{Stealth},thick]
    (0,0) -- (0,4.8) node[above] {FG};

\draw[gray,dashed] (5.2,0) -- (5.2,4.6);
\draw[gray,dashed] (0,2.2) -- (10.7,2.2);

\node[profile] at (2.3,3.5) {
    \textbf{Architectural}\\
    High FG, low FF\\
    Emphasis on appearance
};

\node[profile] at (8,3.5) {
    \textbf{Dual}\\
    High FG, high FF\\
    Visual realism and function
};

\node[profile] at (2.3,.85) {
    \textbf{Non-developed}\\
    Low FG, low FF
};

\node[profile,text=competencygreen] at (8,.85) {
    \textbf{Functional}\\
    Low FG, high FF\\
    Function without literal-copy belief
};

% Architectural to dual: increasing FF
\draw[pathway] (4.15,3.5) -- (6.15,3.5)
    node[midway,above=5pt,explanation,text width=2.6cm] {
        H1a: Visualization makes\\
        functions accessible
    };

% Dual to functional: decreasing FG
\draw[pathway] (10,2.7) -- (10,1.3);

\node[
    explanation,
    anchor=west,
    text width=2.6cm
] at (10.1,2) {
    H1b: Contrasting representations
    challenges literal interpretations
};

\end{tikzpicture}

% \source{Profiles after \citep{ubben2022}.
% Arrows indicate hypothesized pathways, not established effects.}

\note{
The first pathway concerns learners with high FG and low FF.
The visualization may make functional relationships accessible,
for example through wave-plate rotations and polarizer
projections. FF may increase while FG initially remains high,
corresponding to movement from the architectural toward
the dual profile.

The second pathway concerns learners with high FG and high FF.
Contrasting the spatial wave representation with the abstract
sphere may challenge literal interpretations. FG may decrease
while FF remains high, corresponding to movement toward the
functional profile.

These pathways need not occur sequentially or in every learner.
The study measures continuous FF and FG scores rather than
fixed categories. Expert alignment, assessed through FA and GA,
is also part of model understanding and may change along either
pathway, but is not represented by these two axes.

Think-aloud observations may capture spontaneous reasoning
consistent with these pathways during existing applet tasks.
No additional questions or activities are introduced.
}
\end{frame}

\begin{frame}{Research questions}
\footnotesize
\begin{myalert}[Main Research Question 1]
    (How) does the Poincaré sphere visualization change thinking in Fidelity of Gestalt (FG) and/or Functional Fidelity (FF) internal models on polarization?
\end{myalert}
\begin{myalert}[Main Research Question 2]
    Does the Poincaré sphere visualization increase alignment with expert views on Gestalt and Function?
\end{myalert}
\begin{myalert}[Main reseaerch question 3]
To which extend is are the interactions mediated by cognitive load, spatial reasoning ability, and baseline FG?
\end{myalert}


% \keyidea{The sphere may support both functional reasoning and reflection on visual fidelity.}
\note{These are candidate pathways, not established effects or fixed learner trajectories. A wave plate rotates the normalized Stokes vector. An ideal polarizer projects the Jones vector onto its transmission state, followed by normalization for a sphere point. It is not generally a Euclidean projection of a sphere point. The sphere alone does not encode transmission intensity. At zero transmitted intensity, no normalized output state is defined. Cognitive dissonance is one possible explanation for spontaneously expressed conflict, not an inference from confusion alone. Functional understanding can also develop without dissonance.}
\end{frame}
\begin{frame}{Hypotheses}
\footnotesize
\begin{mynote}[H1a: including the architectural profile]
For learners with high FG and low FF, the sphere may make functions accessible: wave plates act as \textbf{rotations} of the state, while polarizers perform \textbf{projections} of the Jones state.
\end{mynote}

\begin{mynote}[H1b: including the dual profile]
For learners with high FG and high FF, contrasting the spatial wave representation with the abstract sphere may challenge literal interpretations and reduce FG while FF stays high.
\end{mynote}
\begin{mynote}[H2: Multimedia effect]
The support with an additional visual representation may increase expert alignment \citep{Rexigel2024}.
\end{mynote}
% \begin{tabularx}{\textwidth}{@{}P{.11\textwidth}P{.30\textwidth}Y@{}}
% \toprule
% & \textbf{Outcome} & \textbf{Sphere versus control, adjusted for baseline}\\
% \midrule
% H1a & Functional Fidelity & Higher FF at post-test.\\
% H1b & Fidelity of Gestalt & Lower FG at post-test.\\
% H2 & Expert alignment & Higher FA and GA (structural/Gestalt alignment).\\
% \bottomrule
% \end{tabularx}\par


% \textbf{Exploratory extension:} What spontaneous reasoning during existing applet tasks helps explain learning and difficulty?
\note{H1/H2 follow the registered directional predictions. The qualitative mechanism analysis is exploratory. Model competencies are theoretical context, not measured outcomes. A pre/post group effect cannot establish mediation by cognitive dissonance.}
\end{frame}


\subsection{Two-phase study}
\begin{frame}{Study structure: pilot, then main study}
\centering
\resizebox{!}{5.35cm}{\begin{tikzpicture}[node distance=7mm,studybox/.append style={font=\fontsize{9}{11}\selectfont,execute at begin node={\hyphenpenalty=10000}}]
\node[studybox,text width=3.1cm] (expert) at (0,1.5) {\textbf{Expert pilot}\\Items, accuracy,\\agreement};
\node[studybox,text width=3.1cm] (novice) at (0,-.5) {\textbf{Novice pilot}\\Usability and\\interpretation};
\node[studybox,text width=3.1cm] (revise) at (0,-2.3) {\textbf{Refine and finalize}\\Instruments and applet};
\draw[flow] (expert.east)--++(.35,0)|-(revise.east);
\draw[flow] (novice)--(revise);
\node[studybox,text width=2.0cm] (pre) at (4.25,1.5) {\textbf{Pre-test}\\FF, FG, FA, GA};
\node[studybox,text width=2.0cm] (random) at (4.25,-.5) {\textbf{Randomize}\\Two conditions};
\node[studybox,text width=2.55cm] (sphere) at (7.4,.35) {Jones + spatial\\\textbf{+ sphere}};
\node[studybox,text width=2.55cm] (control) at (7.4,-1.05) {Jones + spatial\\\textbf{Control}};
\node[studybox,text width=2.0cm] (post) at (10.55,-.35) {\textbf{Post-test}\\Outcomes and other measures};
\draw[flow] (revise.east)--++(.6,0)|-(pre.west);
\draw[flow] (pre)--(random);
\draw[flow] (random.east)--++(.25,0)|-(sphere.west);
\draw[flow] (random.east)--++(.25,0)|-(control.west);
\draw[flow] (sphere.east)--++(.25,0)|-(post.west);
\draw[flow] (control.east)--++(.25,0)|-(post.west);
\node[studybox,draw=competencygreen,text width=4.5cm] (interview) at (5.9,-2.8) {Accompanying think-aloud observations};
\node[studybox,text width=3.0cm] (integrate) at (10.55,-2.8) {Joint interpretation};
% \draw[flow,dashed,draw=competencygreen] (random.south)|-(interview.west);
\draw[flow,draw=competencygreen] (interview)--(integrate);
\draw[flow] (post.south)--(integrate.north);
\node[anchor=west,font=\bfseries\large,text=Accent] at (-1.8,2.65) {Phase 1: Pilot};
\node[anchor=west,font=\bfseries\large,text=Accent] at (3.05,2.65) {Phase 2: Main study};
\draw[gray!50,dashed] (2.25,2.8)--(2.25,-3.5);
\end{tikzpicture}
}
% \source{Accompanying think-aloud observations capture reasoning during existing applet tasks.}
\end{frame}

\begin{frame}{Phase 1: calibrate the instruments and the applet}
\begin{columns}[T,onlytextwidth]
\column{.47\textwidth}
\textbf{Expert pilot}
\begin{itemize}
\item Classify items as FF, FG, both, or neither.
\item Judge factual correctness and the usefulness of analogies.
\item Evaluate wording and agreement.
\end{itemize}
\column{.47\textwidth}
\textbf{Novice pilot}
\begin{itemize}
\item Identify confusing controls and instructions.
\item Examine interpretation of the representations.
\item Assess cognitive load \citep{krieglstein2023} and technological acceptance \citep{abdullah2016}.
\end{itemize}
\end{columns}
\keyidea{Refine questionnaire items, applet instructions, and interaction design before the main comparison.}
\note{The preregistration drops items with IQR greater than 1 and also assesses overall inter-rater reliability. Pilot sizes and the reliability statistic are not specified. It describes think-aloud interviews or questionnaires in the novice pilot. Accompanying main-study observations follow the user's updated plan: existing tasks without additional questions or activities.}
\end{frame}

\begin{frame}{Phase 2: randomized pre/post comparison}
\begin{tabularx}{\textwidth}{@{}P{.18\textwidth}Y@{}}
\toprule
\textbf{Pre-test} & FF, FG, FA, and GA\\
\textbf{Intervention} & Random assignment to guided polarization tasks with or without the added sphere\\
\textbf{Post-test} & FF, FG, FA, GA; cognitive load and spatial reasoning \citep{young2025}, demographic information\\
\bottomrule
\end{tabularx}\par
\vspace{3mm}
\textbf{Population:} second- and third-year physics / related STEM undergraduates.\\[2mm]
\textbf{Target:} $N=218$ recruited; minimum analyzable sample $N=196$.
\note{Recruitment: Leipzig University and affiliated networks. The preregistration specifies single blinding: participants are naive to the manipulation and told broadly that the study concerns visual learning aids in physics. Random assignment occurs when the web app initializes. Registered power assumptions: f-squared .0625, power .80, alpha .05, four outcomes, two groups, two covariates. This slide does not recompute power or set an observation sample size.}
\end{frame}


\begin{frame}{Think-aloud observations: investigating mechanisms}
Students verbalize their thoughts while working through the existing applet tasks.
\begin{itemize}
\item Observe how students attempt to incorporate the Poincar\'e sphere into their thinking.
\item Trace difficulties in connecting graph, symbolic, and sphere representations.
\item Examine spontaneously expressed conflicts between literal interpretations and functional relationships.
\item Follow how these conflicts are resolved, maintained, or bypassed during task engagement.
\end{itemize}
\keyidea{The observations can provide evidence for mechanisms underlying learning and difficulties with mental models.}
\note{The observations may support causal explanations of individual learning episodes. For example, a student might express a literal interpretation, encounter a conflicting representation, and then revise their interpretation. Evidence of confusion alone does not establish cognitive dissonance or identify a dual-type profile. Interpret verbalizations alongside task actions and baseline FF/FG measurements. Students may not verbalize every relevant process. No additional reflection prompts or follow-up questions are planned.}
\end{frame}
% \begin{frame}{Bring quantitative outcomes and verbal reasoning together}
% \begin{columns}[T,onlytextwidth]
% \column{.47\textwidth}
% \textbf{Quantitative strand}
% \begin{itemize}
% \item Estimate condition effects on post-test FF, FG, FA, and GA.
% \item Adjust for baseline scores and spatial reasoning.
% \item Explore initial FG, spatial reasoning, and cognitive load.
% \end{itemize}
% \column{.47\textwidth}
% \textbf{Qualitative strand}
% \begin{itemize}
% \item Trace interpretation, comparison, and revision episodes.
% \item Examine spontaneously expressed realism-related tension.
% \item Examine alternative explanations and contrasting cases.
% \end{itemize}
% \end{columns}
% \keyidea{The experiment estimates condition effects. Observations help explain individual learning episodes.}
% \note{Planned primary approach: MANCOVA. H3 concerns load and spatial reasoning in a tracking subset; H4 concerns initial FG by condition. Interviews offer process evidence, not by themselves causal mediation estimates or a validated competency outcome.}
% \end{frame}
